\section{Abstract}\label{sec:abstract}

People who run open-weight models at home or in a small lab usually end up with a patchwork of servers, proxies, VPNs, and alerts.
Each machine is its own island. Clients must know which box holds which model, a crash in the middle of a long answer
loses the answer, and every open port is attack surface.

Ymir solves this by turning the machines you already own (a GPU workstation, a gaming PC, a Mac, a DGX
Spark-class box, or a CPU-only server) into one private model cloud with a simple single-software \footnoteref{footnote:setup} setup. It starts on a single
machine with one command and grows by linking more of your own devices over Weirdnet, a private,
peer-to-peer network keyed to your identity. The pool behaves as one endpoint: requests go to a
node that has the model ready, models are spread across the pool, and only devices linked to the owner's identity can use or command the pool.

\begin{table}[htbp]
\centering\small
\begin{tabularx}{\linewidth}{@{}lY@{}}
\toprule
Repository & \url{https://github.com/WeirdBrains/ymir} \\
License & BSD-3-Clause, copyright \WB \\
Language & Rust (edition 2024, Rust 1.99) \\
Binaries & \code{ymird} (daemon), \code{ymirctl} (CLI) \\
Client & Bragi (separate repo, v0.1.0 released) \\
Release state & No tagged Ymir release yet; build from source \\
\bottomrule
\end{tabularx}
\end{table}

\footnote{\label{footnote:setup}You can add additional engines (e.g. vLLM, llama-cpp, etc.) and Ymir will automatically use them.}
